Congressional redistricting occurs every decade, yet research has focused almost exclusively on single-census optimization. This temporal myopia ignores a critical question: as demographics shift over ten years, how stable do district boundaries remain? We address this gap through the first systematic comparison of graph partitioning methods across census cycles.

Our analysis of five southern states over 2010-2020 reveals a fundamental finding: \textbf{hierarchical recursive bisection provides substantially greater temporal stability than flat n-way partitioning.} Recursive methods maintain 80\% tract retention versus 70\% for n-way (Cohen's $d=1.2$, $p<0.001$), reduce population disruption by 10 percentage points, and preserve 94\% of top-level geographic splits. Communities of interest experience 14 percentage points less disruption.

This stability advantage stems from hierarchical structure preservation. Recursive bisection's binary tree creates nested regions reflecting fundamental geography (urban-rural divides, historical regions, natural boundaries). These top-level splits persist across decades because they encode geographic constants. When demographics shift, adjustments occur within regions without cascading across all boundaries. By contrast, n-way's flat structure re-optimizes all districts simultaneously, causing global disruption.

The implications are both practical and theoretical. \textbf{Practically}, jurisdictions face a performance-stability tradeoff: n-way offers +4.3 percentage points better 2020 minority concentration but -10 to -13 points worse 2020-2030 stability. For comfortable VRA margins, recursive bisection's stability benefits outweigh marginal performance costs. For borderline cases, n-way's immediate gains may dominate. Hybrid approaches—initializing n-way with recursive structure—could balance both objectives.

\textbf{Theoretically}, our findings generalize beyond redistricting. Any domain requiring periodic re-partitioning (school districts, service territories, sales regions) faces similar temporal dynamics. When re-partitioning frequency is low relative to timescale, stability gains value. Hierarchical methods provide continuity through structural preservation, while flat methods sacrifice disruption for per-instance optimization. This constitutes a fundamental tradeoff in dynamic partitioning problems.

Our work generates testable predictions for 2030 redistricting: states using recursive methods in 2020 should experience less 2030 disruption, and preserving 2020 tree structures should improve stability. When 2030 census data arrives, researchers can validate these predictions, testing our mechanistic explanation of hierarchical structure advantages.

Future research should extend temporally (adding 2000 data for three-cycle analysis), geographically (all 50 states), and algorithmically (hybrid methods, stability-aware optimization). Enhanced metrics—representative-constituent continuity, voter awareness surveys, administrative cost models—would capture socio-political dimensions beyond geographic stability. Integrating demographic projections could enable proactive stability optimization, jointly planning 2020 and predicted 2030 boundaries.

\textbf{In sum}: redistricting method choice should consider not only current optimization but also future continuity. Recursive bisection's hierarchical structure provides substantial temporal stability advantages—advantages likely to grow in importance as redistricting reform efforts emphasize community preservation and administrative efficiency alongside traditional criteria. When communities of interest matter, when administrative disruption imposes costs, and when long-term planning horizons extend beyond single census cycles, hierarchical methods merit serious consideration.

The question is no longer just ``which method produces better districts today?'' but ``which method produces districts that remain coherent tomorrow?'' Our findings suggest that for temporal stability, hierarchical structure makes all the difference.
